\documentclass[a4paper,12pt]{article}

\usepackage[left=2cm, right=2cm, top=3cm, bottom=3cm]{geometry} %Para poner los márgenes


\usepackage{float} %Para hacer que la megatabla no se vaya volando a la última página y forzar que se quede usando [H]

% Idioma español y codificación
\usepackage[spanish]{babel} % Traduce palabras como "Contenido", "Capítulo", etc.
\usepackage[utf8]{inputenc} % Permite escribir caracteres en español (á, é, ñ, etc.)
\usepackage[T1]{fontenc}    % Mejora la salida de caracteres acentuados en PDF

% Paquetes matemáticos
\usepackage{amsmath}  % Permite usar comandos avanzados de matemáticas
\usepackage{amssymb}  % Agrega símbolos matemáticos adicionales
\usepackage{amsthm}   % Permite definir teoremas, lemas, etc.
\usepackage{mathrsfs} % Fuente adicional para letras matemáticas

% Para mejor tipografía en matemáticas
\usepackage{mathtools}

% Si necesitas gráficos
\usepackage{graphicx} 

% Si usas referencias cruzadas en el documento
\usepackage[hidelinks]{hyperref}  % Enlaces clicables en PDF

%Para poner el color de los enlaces azul
\hypersetup{
    colorlinks=true,       % Habilita los enlaces de colores
    linkcolor=black,        % Color de los enlaces internos (por ejemplo, referencias a secciones)
    urlcolor=magenta,      % Color de los enlaces a direcciones web
    citecolor=red,         % Color de las referencias bibliográficas
    filecolor=cyan,        % Color de los enlaces a archivos
}

\renewcommand{\contentsname}{Índice} % Cambia "Contenido" por "Índice"

\usepackage{listings}
\usepackage{xcolor} % Para definir colores, opcional

% Configuración básica de listings:
\lstset{
    language=C,                    % Especifica el lenguaje (puede ser Python, Java, etc.)
    basicstyle=\ttfamily\small,     % Fuente de tipo monoespaciada
    numbers=left,                  % Números de línea a la izquierda
    numberstyle=\tiny,             % Estilo de los números de línea
    stepnumber=1,                  % Mostrar número en cada línea
    numbersep=10pt,                % Espacio entre número y código
    backgroundcolor=\color{lightgray}, % Color de fondo
    showspaces=false,              % No resaltar espacios
    showstringspaces=false,        % No resaltar espacios en cadenas
    showtabs=false,                % No resaltar tabulaciones
    frame=single,                  % Agregar un marco alrededor del código
    rulecolor=\color{black},       % Color del marco
    captionpos=b,                  % Posición de la leyenda (b de bottom)
    breaklines=true,               % Saltos automáticos de línea
    breakatwhitespace=false,
    title=\lstname,                % Muestra el nombre del archivo, si se especifica
    keywordstyle=\color{blue},     % Color de las palabras clave
    commentstyle=\color{purple},    % Color de los comentarios
    stringstyle=\color{red}        % Color de las cadenas
}

\usepackage{caption}
\captionsetup[table]{labelformat=empty} %Para que no salgan numeradas las tablas
\captionsetup[lstlisting]{labelformat=empty} %Para que no salgan numeradas los fragmentos de código
\captionsetup[table]{labelformat=empty} %Para que no salgan numeradas las tablas
\captionsetup[figure]{labelformat=empty} %Para que no salgan numeradas las gráficas
\setlength{\parindent}{0.75cm} %Para que la sangría no sea tan grande
 % Todos los preámbulos

% =====================
% Índice
% =====================
\begin{document}
\begin{titlepage}
	\centering

	\includegraphics[width=0.3\textwidth]{Imagenes/logo_ugr.png} % Ajusta el tamaño según necesites (0.3\textwidth es solo un ejemplo)
	\vspace{1cm} % Ajusta el espacio según lo necesites

	{\bfseries\LARGE Universidad de Granada \par}
	\vspace{1cm}
	{\scshape\Large Facultad de Ingeniería Informática y telecomunicaciones \par}
	\vspace{3cm}
	{\scshape\Huge Práctica 1: Eficiencia de algoritmos \par}
	\vspace{3cm}
	{\itshape\Large Doble Grado Ingeniería Informática y Matemáticas \par}
	\vfill
	{\Large Autores: \par}
	{\Large Adolfo Martínez Olmedo, Pablo Delgado Galera, Marcos Baena Solar \par}
	\vfill
	{\Large Marzo 2025 \par}
\end{titlepage}
 % Archivo con la portada

% =====================
\section{Introducción}
% =====================
En esta práctica se desarrolla la resolución de tres problemas diferentes: 

\begin{itemize}
	\item Menor número
	\item Menor distancia
	\item Menor envolvente convexa
\end{itemize}

Lo haremos de dos maneras distintas, primero haremos una aproximación mediante algoritmos 
de fuerza bruta y después aplicaremos lo aprendido en clase, desarrollando soluciones basadas 
en divide y vencerás.


% =====================
\section{Objetivos}
% =====================
Los objetivos de esta práctica son los siguientes:
\begin{itemize}
	\item Comprender la técnica de Divide y Vencerás y sus ventajas.
	\item Comparar con la estrategia de fuerza bruta y analizar la complejidad.
	\item Implementar ambos enfoques (fuerza bruta y Divide y Vencerás) para cada problema.
	\item Experimentar con el umbral de la técnica de Divide y Vencerás.
\end{itemize}

% =====================
\section{El número más pequeño}
\label{sec:El número más pequeño}
% =====================
\subsection{Análisis de Fuerza Bruta}

Nuestra idea para desarrollar el algoritmo de fuerza bruta se basa en ir recorriendo el vector de puntos tantas veces como cifras tenga el número que busquemos, y en cada recorrido se busca el número más pequeño, se elimina y se procede a hacer lo mismo en la siguiente iteración.

El código es el siguiente:

\begin{lstlisting}[language=C, caption={Brute Force de \textbf{"El número más pequeño"}}]
	int getSmallestNumber(vector<int> v, int k){
		int smallestNumber = std::numeric_limits<int>::max();
		for (int i = 0; i < v.size()-1; ++i){
			string num;
			num.append(to_string(v[i]));
			for (int j = i+1; j-i <= k  && v.size()-1 - j >= k; ++j){
				num.append(to_string(v[j]));
			}
			if(stoi(num) < smallestNumber)
				smallestNumber = stoi(num);   
		}
	
		return smallestNumber;
	}	
	 \end{lstlisting}

Tras el análisis teórico comprobamos que este código tiene una eficiencia de $O(k\cdot n)$.
\subsection{Análisis de Divide y Vencerás}
Nuestra idea para desarrollar el algoritmo de Divide y Vencerás se basa en orderar el vector utilizando quickSort (para dividir y vencer) y una vez lo tenemos ordenado simplemente cogemos los k primeros números.

El código es el siguiente:

\begin{lstlisting}[language=C, caption={Divide y vencerás de \textbf{"El número más pequeño"}}]
void quicksort(vector<int> &v, int left, int right)
{
	if (left >= right)
		return;
	int pivot = v[right];
	int partitionIndex = left;
	for (int i = left; i < right; ++i)
	{
		if (v[i] < pivot)
		{
			swap(v[i], v[partitionIndex]);
			++partitionIndex;
		}
	}
	swap(v[partitionIndex], v[right]);
	quicksort(v, left, partitionIndex - 1);
	quicksort(v, partitionIndex + 1, right);
}

int getSmallestNumber(vector<int> v, int k)
{
	quicksort(v, 0, v.size() - 1);
	string result = "";
	for (int i = 0; i < k; ++i)
	{
		result += to_string(v[i]);
	}
	return stoi(result);
}	
	 \end{lstlisting}

Tras el análisis teórico comprobamos que este código tiene una eficiencia de $O(n\cdot \log (n))$

\subsection{Configuración de las Pruebas}
\begin{itemize}
	\item Descripción del entorno de ejecución (CPU, memoria, compilador, etc.).
	\item Conjunto de datos utilizados para las pruebas (tamaño, forma de generarlos).
\end{itemize}

\subsection{Tablas y Gráficos de Rendimiento}
Insertar aquí las tablas y/o gráficos que muestren los tiempos de ejecución,
uso de memoria, etc.

\subsection{Análisis de los Resultados}
Tras el analisis teórico y visualizar la gráfica comparativa, vemos que el algoritmo de Fuerza 
Bruta consigue resolver el problema de forma más eficiente que el algoritmo Divide y Vencerás. 
Esto se debe a que el primer algortimo depende de una constante k que consideramos siempre menor 
que n, y por tanto no afecta a la eficiencia. Mientras que en el segundo algoritmo podemos apreciar 
que quicksort tiene la eficiencia del peor caso, ya que al solo usar dígitos del 1 al 9, encontramos
en el vector una gran cantidad de números repetidos lo cual afecta gravemente a la eficiencia de este 
algortimo de ordenación.

\begin{itemize}
	\item Comparación cualitativa (efectividad, facilidad de implementación).
	\item Comparación cuantitativa (tiempos de ejecución, consumo de memoria).
	\item Conclusiones sobre el umbral experimental.
\end{itemize}



% =====================
\section{El par de puntos más cercano}
% =====================
\subsection{Análisis de Fuerza Bruta}
Nuestra idea para desarrollar el algoritmo de fuerza bruta se basa en un doble bucle for el cual hace todas las distancias de todos los puntos con todos, hasta encontrar la mínima distancia. 

El código es el siguiente:

\begin{lstlisting}[language=C, caption={Brute Force de \textbf{"El par de puntos más cercano"}}]
struct Point {
	double x, y;
};

double dist(Point p1, Point p2) {
	return sqrt( (p1.x - p2.x)*(p1.x - p2.x) + (p1.y - p2.y)*(p1.y - p2.y) );
}


double bruteForce(const vector<Point>& P, int l, int r) {
	double minDist = DBL_MAX;
	for (int i = l; i <= r; i++) {
		for (int j = i+1; j <= r; j++) {
			double d = dist(P[i], P[j]);
			if (d < minDist) {
				minDist = d;
			}
		}
	}
	return minDist;
}
\end{lstlisting}

Tras el análisis teórico comprobamos que este código tiene una eficiencia de $O(n^2)$.

\subsection{Análisis de Divide y Vencerás}
Nuestra idea para desarrollar el algoritmo de Divide y Vencerás se basa en (Pablo Indio complétalo).

El código es el siguiente:

\begin{lstlisting}[language=C, caption={Divide y vencerás de \textbf{"El par de puntos más cercano"}}]
	void swap(Point* a, Point* b) {
		Point t = *a;
		*a = *b;
		*b = t;
	}
	
	int partition(vector<Point>& arr, int low, int high) {
		double pivot = arr[high].x;
		int i = (low - 1);
	
		for (int j = low; j <= high - 1; j++) {
			if (arr[j].x < pivot) {
				i++;
				swap(&arr[i], &arr[j]);
			}
		}
		swap(&arr[i + 1], &arr[high]);
		return (i + 1);
	}
	
	void quickSort(vector<Point>& arr, int low, int high) {
		if (low < high) {
			int pi = partition(arr, low, high);
	
			quickSort(arr, low, pi - 1);
			quickSort(arr, pi + 1, high);
		}
	}
	
	void sort_by_x(vector<Point>& points) {
		quickSort(points, 0, points.size() - 1);
	}
	
	double bruteForce(const vector<Point>& P, int l, int r) {
		double minDist = DBL_MAX;
		for (int i = l; i <= r; i++) {
			for (int j = i+1; j <= r; j++) {
				double d = dist(P[i], P[j]);
				if (d < minDist) {
					minDist = d;
				}
			}
		}
		return minDist;
	}
	
	
	double closestUtil(vector<Point> &PX, int l, int r) {
		// Si el subrango es peque, fuerza bruta
		if (r - l <= 3) {
			return bruteForce(PX, l, r);
		}
	
		int mid = (l + r) / 2;
		double midX = PX[mid].x;
	
		// bestPair se va actualizando; 
		// para tener la minima distancia, llamo recursivamente a izq y der
		double dl = closestUtil(PX, l, mid);
		double dr = closestUtil(PX, mid+1, r);
	
		double d = dl < dr ? dl : dr;
		// Construir el strip (puntos que quedan en [midX - d, midX + d])
		vector<Point> strip;
		strip.reserve(r - l + 1);
		for (int i = l; i <= r; i++) {
			if (fabs(PX[i].x - midX) < d) {
				strip.push_back(PX[i]);
			}
		}
		// Ordenar la franja por Y
		sort(strip.begin(), strip.end(), [](auto &a, auto &b){
			return a.y < b.y;
		});
	
		// Revisar cada punto del strip con sus siguientes
		// mientras la diferencia en Y sea < d
		for (int i = 0; i < (int)strip.size(); i++) {
			for (int j = i+1; j < (int)strip.size() && (strip[j].y - strip[i].y) < d; j++) {
				double distStrip = dist(strip[i], strip[j]);
				if (distStrip < d) {
					d = distStrip;
				}
			}
		}
	
		return d;
	}
\end{lstlisting}

Tras el análisis teórico comprobamos que este código tiene una eficiencia de (PABLO INDIO HAZLO)

\subsection{Detalles de Implementación}

\subsection{Configuración de las Pruebas}
\begin{itemize}
	\item Descripción del entorno de ejecución (CPU, memoria, compilador, etc.).
	\item Conjunto de datos utilizados para las pruebas (tamaño, forma de generarlos).
\end{itemize}

\subsection{Tablas y Gráficos de Rendimiento}
Insertar aquí las tablas y/o gráficos que muestren los tiempos de ejecución,
uso de memoria, etc.

\subsection{Análisis de los Resultados}
\begin{itemize}
	\item Comparación cualitativa (efectividad, facilidad de implementación).
	\item Comparación cuantitativa (tiempos de ejecución, consumo de memoria).
	\item Conclusiones sobre el umbral experimental.
\end{itemize}


% =====================
\section{La envolvente convexa}
% =====================
\subsection{Análisis de Fuerza Bruta}

Nuestra idea para desarrollar el algoritmo de fuerza bruta se basa en empezar seleccionando el punto más a la izquierda, el cual sabes que siempre va a estar en la envolvente convexa, seguidamente seleccionas otro punto y compruebas si la recta que une los dos puntos deja el resto de puntos a la izquierda, una vez encuentras ese punto, es nuevo punto actúa como el primero que seleccionaste, y haces esto hasta que encuentres todos los puntos pertenecientes a la envolvente convexa.

El código es el siguiente:

\begin{lstlisting}[language=C, caption={Brute Force de \textbf{"La envolvente convexa"}}]
	// Estructura para representar un punto
	struct Point {
		float x, y;
	};
	
	typedef pair<int, int> ParPuntos;
	
	// Funcion que calcula el convex hull por fuerza bruta
	set<ParPuntos> convexHullBruteForce(const vector<Point>& points) {
		int n = points.size();
		set<ParPuntos> hull;
		
		// Recorremos cada par de puntos
		for (int i = 0; i < n; i++) {
			for (int j = i + 1; j < n; j++) {
				int left = 0, right = 0;
				
				// Verificamos la posicion de cada otro punto respecto a la linea formada por points[i] y points[j]
				for (int k = 0; k < n; k++) {
					if (k == i || k == j)
						continue;
						
					float position = (points[j].x - points[i].x) * (points[k].y - points[i].y) -
								   (points[j].y - points[i].y) * (points[k].x - points[i].x);
					
					if (position > 0)
						left++;
					else if (position < 0)
						right++;
				}
				
				// Si todos los puntos se encuentran de un mismo lado (o son colineales), se adjuntan los puntos de la arista
				if (left == 0 || right == 0) {
					hull.insert({points[i].x, points[i].y});
					hull.insert({points[j].x, points[j].y});
				}
			}
		}
		
		return hull;
	}
\end{lstlisting}

Tras el análisis teórico comprobamos que este código tiene una eficiencia de $O(n^3)$

\subsection{Análisis de Divide y Vencerás}
Nuestra idea para desarrollar el algoritmo de Divide y Vencerás se basa en (pablo भारतीय इसे पूरा करें)
\subsection{Detalles de Implementación}

\subsection{Configuración de las Pruebas}
\begin{itemize}
	\item Descripción del entorno de ejecución (CPU, memoria, compilador, etc.).
	\item Conjunto de datos utilizados para las pruebas (tamaño, forma de generarlos).
\end{itemize}

\subsection{Tablas y Gráficos de Rendimiento}
Insertar aquí las tablas y/o gráficos que muestren los tiempos de ejecución,
uso de memoria, etc.

\subsection{Análisis de los Resultados}
\begin{itemize}
	\item Comparación cualitativa (efectividad, facilidad de implementación).
	\item Comparación cuantitativa (tiempos de ejecución, consumo de memoria).
	\item Conclusiones sobre el umbral experimental.
\end{itemize}

En esta sección se describen los diferentes problemas que se abordarán.

% =====================
\section{Conclusiones}
% =====================
Tras esta práctica hemos descubierto que la estrategia de desarrollo de algoritmos Divide y Vencerás es útil en un amplio repertorio de casos pero esto no quiere decir que siempre sea mejor que un algoritmo fuerza bruta, esto lo podemos comprobar en \ref{sec:El número más pequeño}. Aún así en la gran mayoría de casos aplicar Divide y Vencerás 
supone una mejoría en la eficiencia del código.

\end{document}

